\section{A worked dependent connection}
\label{app:worked}

Let the roles be $A,B,C$ and use the script in (1), with
\[
 R_A(n,v,j)=\Bool,\qquad R_B(n,v,j)=\mathsf{Fin}(n),
 \qquad R_C(n,v,j)=\mathsf{Vec}(\Bool,n).
\]
Its local types are
\begin{align*}
 \Local{A}{G}
   &=\sum_{n:\Nat_{>0}}\prod_{v:\mathsf{Vec}(\Bool,n)}
                    \prod_{j:\mathsf{Fin}(n)}\Bool,\\
 \Local{B}{G}
   &=\prod_{n:\Nat_{>0}}\sum_{v:\mathsf{Vec}(\Bool,n)}
                    \prod_{j:\mathsf{Fin}(n)}\mathsf{Fin}(n),\\
 \Local{C}{G}
   &=\prod_{n:\Nat_{>0}}\prod_{v:\mathsf{Vec}(\Bool,n)}
                    \sum_{j:\mathsf{Fin}(n)}\mathsf{Vec}(\Bool,n).
\end{align*}
Choose implementations
\begin{align*}
 p_A&=(2,\lambda v.\lambda j.\mathsf{index}(v,j)),\\
 p_B&=\lambda n.(\mathsf{replicate}(n,\mathsf{true}),\lambda j.j),\\
 p_C&=\lambda n.\lambda v.(0_n,v),
\end{align*}
where $0_n:\mathsf{Fin}(n)$ exists because $n$ is positive. The first
interaction selects $2$ and instantiates both receivers at that number.
The next selects Bob's vector $(\mathsf{true},\mathsf{true})$. The third
selects Carol's $0_2$. The final return record is
\[
 \bigl(\mathsf{true},0_2,
       (\mathsf{true},\mathsf{true})\bigr),
\]
indexed by the common trace
$(2,(\mathsf{true},\mathsf{true}),0_2)$.

Neither receiver was required to guess a dependent type from a value it had
not received. The first substitution determines the vector and index types;
the next two substitutions determine the return families. The example
exercises three different owners and all the clauses of the interaction rule.

\section{Two distinctions in finite models}

\paragraph{Policy tables and plans.}
For three binary choices there are $2^{1+2+4}=128$ complete action-policy
tables for any fixed owner assignment: one row at the first occurrence, two
at the second, and four at the third. Distributing the three occurrences over
three roles gives $3^3=27$ assignments. The accompanying check compares all
$3,456$ full profiles with their reduced plans and tests each of the six role
permutations. Terminal returns are fixed pure functions of the trace in this
finite experiment. The comparison is of complete evaluations, as required by
\cref{thm:adequacy}, not an assertion that the table and plan carriers have
the same cardinality.

\paragraph{Everyone knows and common knowledge.}
It is insufficient that all players know a fact at the actual history. They
may consider histories possible at which another player does not know it.
Condition (25) rules this out by preserving the public readout at every
indistinguishability edge. A chain may alternate arbitrarily among the roles,
yet cannot change that readout. This is why the proof establishes common
knowledge rather than just the first level of mutual knowledge.

\section{Result inventory}

\begin{center}
\small
\begin{tabular}{>{\raggedright\arraybackslash}p{0.35\linewidth}>{\raggedright\arraybackslash}p{0.56\linewidth}}
\toprule
Result & Scope and assumptions\\
\midrule
Substitution, \cref{lem:substitution} & Finite public scripts over a data
calculus with ordinary dependent substitution.\\[3pt]
Typed interaction, \cref{thm:interaction,cor:preservation} & A complete
profile implementing one common public script.\\[3pt]
Termination and uniqueness, \cref{thm:termination} & Closed complete
connections; strong normalization, confluence, and canonical forms in the
pure data calculus.\\[3pt]
Role symmetry, \cref{thm:equivariance} & Bijective renaming of roles and the
corresponding endpoint and result transport.\\[3pt]
Boundary sequencing, \cref{prop:sequence,thm:associativity} & An explicit
public return boundary, retaining the complete intermediate record.\\[3pt]
Definability, \cref{prop:definability} & The finite public interpretation and
its execution macros; original data judgments.\\[3pt]
Phase preservation and binding,
\cref{lem:phase-preservation,prop:binding} & Well-formed phase contexts,
typed pure implementations, and fresh immutable cell names.\\[3pt]
Hiding, \cref{thm:hiding} & Two runs differing at a protected binding;
a prefix sealed against direct and dependency-mediated release.\\[3pt]
Common knowledge, \cref{thm:common-knowledge} & A fixed public stage and a
readout factoring through every relevant player's observation.\\[3pt]
Strategy adequacy, \cref{thm:adequacy} & Reduction of full strategies to
endpoint plans; equality of complete execution, not of SPE status.\\[3pt]
Rational algebra, \cref{prop:rational,lem:rational-algebra,cor:no-boolean-product}
& Finite perfect-information one-winner trees, pure full SPE profiles,
and no strategic terminal-return stage.\\
\bottomrule
\end{tabular}
\end{center}

The general open-interface rules (27)--(28), their multiparty cut theorem,
and a dependent typing system for arbitrary private networks are development
targets rather than entries in this inventory.
